% =============================================================================
%  OP2: FEP-as-lens scope — skeleton paper
%  STATUS: UNPROVEN. Hardest of the open problems; depends on what
%          "informative FEP" is committed to mean.
% =============================================================================
\documentclass[11pt]{article}

\usepackage[utf8]{inputenc}
\usepackage[T1]{fontenc}
\usepackage[expansion=false]{microtype}
\usepackage{amsmath, amssymb, amsthm, mathtools}
\usepackage[margin=1in]{geometry}
\usepackage{enumitem}
\usepackage[hidelinks]{hyperref}

\theoremstyle{plain}
\newtheorem{theorem}{Theorem}
\newtheorem*{namedopenproblem}{Open Problem}
\theoremstyle{definition}
\newtheorem{definition}[theorem]{Definition}

\title{The Scope of the Free Energy Principle as a Parameterisation \\
       \large (UNPROVEN --- skeleton paper for OP2)}
\author{Liam McGuire\thanks{Unaffiliated researcher, Seattle, USA. liammcg44@gmail.com}}
\date{\today}

\begin{document}
\maketitle

\begin{abstract}
This paper is a \textbf{skeleton} that records OP2 from the
scaffolding paper~\cite{emergent_systems_2026}. \textbf{No proof is
attempted.} OP2 is technically heavy: the answer requires committing
to a precise meaning of ``informative FEP parameterisation,'' which
the field has not converged on. The result is UNPROVEN.
\end{abstract}

\section{Statement (UNPROVEN)}

\begin{quote}
\emph{``OP2 (FEP-as-lens scope): characterise precisely the
substrate--dynamics pairs $(X, \Phi)$ for which the FEP is
informative as a parameterisation of viability and observer slots.''}
\end{quote}
\noindent (Scaffolding paper~\cite{emergent_systems_2026},
\S\texttt{sec:c1}-onward; the open-problem statement is in the
Conjectures section.)

\section{Plan of attack (no proof attempted)}

Following the proof pattern of Aguilera et
al.~\cite{aguilera2022particular}:

\begin{enumerate}[leftmargin=2em,itemsep=4pt]
  \item \textbf{Tractable class restriction (UNPROVEN).} Restrict to
    weakly-coupled linear / Gaussian-NESS SDEs.
  \item \textbf{Translate FEP requirements to matrix conditions
    (UNPROVEN).} ``MB exists'' becomes a block-zero condition on the
    NESS precision $\Pi$; ``non-trivial variational bound'' becomes
    a non-degeneracy condition on the synchronisation map.
  \item \textbf{Codimension computation (UNPROVEN).} Show the
    simultaneous block-zero + solenoidal constraints cut out a
    measure-zero submanifold (Aguilera) OR that a relaxation
    recovers a positive-measure regime (Sakthivadivel; Heins--Da Costa).
  \item \textbf{Iff statement (UNPROVEN).} Candidate:
    \emph{$(X, \Phi)$ admits an informative FEP parameterisation iff
    (i) sparse-causal-coupling-admits-MB regime, (ii) non-degenerate
    synchronisation manifold, (iii) tractability of the implicit
    generative model.}
\end{enumerate}

\section{Connection to N1}

Under the pivot framing, OP2 is the FEP-specialised case of N1: it
characterises which $\mathcal{O}_W$-presentations admit informative
FEP parameterisations.

\section{Status and follow-up}

\textbf{Current status: UNPROVEN.} The hardest of the open problems.
Requires committing to what ``informative'' means.

\bibliographystyle{plain}
\bibliography{references}

\end{document}
